% SPDX-License-Identifier: CC-BY-SA-4.0
% Part: Lexical Analysis
% Section: Minimal automata
% Challenge exercise: Equivalence between regular expressions

\begingroup

\begin{exercice}[Équivalence entre expressions rationnelles]\label{exo:lexing/minimisation/equivalence}

  \begin{question}
  \item En appliquant l'algorithme de Thompson, donner un automate fini non déterministe reconnaissant \linebreak$\mathcal{S}((a^\star \mid b^\star)^\star)$. 
  \item Déterminiser l'automate obtenu en utilisant l'algorithme de Rabin et Scott. 
  \item Minimiser l'automate obtenu en utilisant l'algorithme de Moore. 
  \item En suivant les mêmes étapes, donner l'automate minimal reconnaissant $\mathcal{S}((a \;|\; b)^\star)$. 
  \item En déduire que $(a^\star \;|\; b^\star)^\star \equiv (a \;|\; b)^\star$. 
  \item Proposer une méthode générale pour décider si deux expressions rationnelles décrivent le même langage,
    en utilisant les algorithmes vus en cours. 
  \end{question}

\end{exercice}

\begin{correction}

  \begin{question}

  \item \begin{tikzpicture}[automaton, baseline=(0.base)]
    \state[initial]   (0)  at (0,3) {$0$}; 
    \state            (1)  at (1,2) {$1$}; 
    \state            (2)  at (2,3) {$2$}; 
    \state            (3)  at (3,3) {$3$}; 
    \state            (4)  at (3,4) {$4$}; 
    \state            (5)  at (4,3) {$5$}; 
    \state            (6)  at (2,1) {$6$}; 
    \state            (7)  at (3,1) {$7$};
    \state            (8)  at (3,0) {$8$};
    \state            (9)  at (4,1) {$9$};
    \state            (10) at (5,2) {$10$};
    \state[accepting] (11) at (0,1) {$11$}; 
    
    \path (0) edge                 node       {$\varepsilon$} (1);
    \path (1) edge                 node       {$\varepsilon$} (11);
    \path (1) edge                 node       {$\varepsilon$} (2);
    \path (1) edge                 node       {$\varepsilon$} (6);
    \path (2) edge                 node       {$\varepsilon$} (3);
    \path (3) edge                 node       {$\varepsilon$} (5);
    \path (6) edge                 node       {$\varepsilon$} (7);
    \path (7) edge                 node       {$\varepsilon$} (9);
    \path (5) edge                 node       {$\varepsilon$} (10);
    \path (9) edge                 node       {$\varepsilon$} (10);
    \path (10)edge                 node       {$\varepsilon$} (1);
    \path (3) edge[bend left=8mm]  node       {$a$} (4);
    \path (4) edge[bend left=8mm]  node       {$\varepsilon$} (3);
    \path (7) edge[bend right=8mm] node[swap] {$b$} (8);
    \path (8) edge[bend right=8mm] node[swap] {$\varepsilon$} (7);
  \end{tikzpicture}

  \item \begin{tikzpicture}[automaton, baseline=(A.base)]
    \state[initial, accepting] (A) at (1,1) {$A$}; 
    \state[accepting]          (B) at (0,0) {$B$}; 
    \state[accepting]          (C) at (2,0) {$C$}; 
    
    \path (A) edge                 node[swap] {$a$} (B);
    \path (A) edge                 node       {$b$} (C);
    \path (C) edge[bend left=3mm]  node       {$a$} (B);
    \path (B) edge[bend left=3mm]  node       {$b$} (C);
    \path (B) edge[loop left]      node       {$a$} (B);
    \path (C) edge[loop right]     node       {$b$} (C);
  \end{tikzpicture}

  \item \begin{tikzpicture}[automaton, baseline=(q.base)]
    \state[initial, accepting] (q) {$q$}; 
    \path                      (q) edge[loop right] node {$a, b$}   (q);
  \end{tikzpicture}

  \item L'automate non-déterministe est le suivant. L'automate déterminisé et l'automate minimal sont les mêmes que dans le cas précédent.
    \begin{tikzpicture}[automaton, baseline=(0.base)]
      \state[initial]   (0)  at (0,2) {$0$}; 
      \state            (1)  at (1,1) {$1$}; 
      \state            (2)  at (2,2) {$2$}; 
      \state            (5)  at (3,2) {$3$}; 
      \state            (6)  at (2,0) {$4$}; 
      \state            (9)  at (3,0) {$5$};
      \state            (10) at (4,1) {$6$};
      \state[accepting] (11) at (0,0) {$7$}; 
      
      \path (0)  edge node {$\varepsilon$} (1);
      \path (1)  edge node {$\varepsilon$} (11);
      \path (1)  edge node {$\varepsilon$} (2);
      \path (1)  edge node {$\varepsilon$} (6);
      \path (5)  edge node {$\varepsilon$} (10);
      \path (9)  edge node {$\varepsilon$} (10);
      \path (10) edge node {$\varepsilon$} (1);
      \path (2)  edge node {$a$} (5);
      \path (6)  edge node {$b$} (9);
    \end{tikzpicture}

  \item Les langages décrits par les deux expressions rationnelles ont le même automate minimal, donc ils sont égaux. 

  \item En général, on utilise le théorème d’unicité de l’automate minimal :
  on calcule les automates finis minimaux associés aux deux expressions rationnelles.
  Elles reconnaissent le même langage si, et seulement si, ces automates minimaux sont isomorphes.

  \end{question}

\end{correction}

\endgroup
\endinput
